% 复制本文件到 chapters/，替换全部占位符，并在 main.tex 中加入 \input。
% 主要资料：书名、版本或资料名称
% 精确范围：章节、印刷页码与 PDF 页码
% 人工核验：否

\section{第 X.Y 节：主题标题}
\label{section:QM-CXX-SYY}

\studymeta
  {QM-CXX-SYY}
  {YYYY-MM-DD}
  {主要资料名称}
  {第 X 章 \S X.Y；印刷页 pp.~XX--YY；PDF pp.~AA--BB}
  {待确认}

\subsection{本节主线}

用自己的话说明这一节解决什么问题，以及它与前后内容的关系。

\notetopic{QM-CXX-SYY-T01}{知识点标题}

\keyidea{用一两句话写出核心思想。}

\begin{definition}
写出精确定义，并明确对象、标量域、符号和假设。
\end{definition}

\subsubsection{动机与机制}

解释为什么需要这一概念，以及推理怎样一步步进行。

\subsubsection{推导补全}

教材若省略推导，在这里补全步骤，并说明这是笔记补充。

\subsubsection{边界与常见误区}

记录结论的适用条件、容易混淆的对象和可能失效的假设。

\subsubsection{相关练习}

\relatedexercise{QM-CXX-SYY-E01}{对应教材练习或自拟检查题}。

\subsection{一分钟回顾}

用三到五句话总结本单元，不引入未经说明的新结论。
